\section{Materials}\label{sec:materials}

\subsection{Five libraries with released MEL slices}

We use the five collections for which the MEL family is released with per-class labels and more
than one slice per recording: CTIS for Chinese instrument recognition, music genre for
popular-song excerpts, pianos for eight piano brands, CNPM for Chinese pentatonic-mode classification, and
bel canto for bel-canto and folk singing. Three of the five --- CTIS, CNPM and bel canto --- are
among the six datasets evaluated in the audited article; \emph{music genre} and \emph{pianos}
belong to the same published collection of libraries but are not covered by it, so contrasts in
those two libraries audit the same kind of split under a framework that does not report on them.
Together they release 110{,}415 slices from 7{,}827 recordings in
252 source classes; folding the duplicate takes described below leaves 7{,}825 canonical
recordings, which is the footing of the library inventory in the supplement. Of these,
109{,}831 slices from 7{,}756 recordings in 234 classes survive the
minimum-recording rule and placeholder removal described in
Section~\ref{sec:capacity}, listed in Table~\ref{tab:panels}; the remainder is excluded under
these study-specific rules. In CTIS the rule leaves 42{,}509 of the 43{,}054 released slices in
201 classes, and four of those slices repeat a content hash already claimed by the sealed panel,
because the same take is released under two technique names; those four are dropped when the
partitions are built, so the four partitions carry 42{,}505. The four candidate near-single-slice collections are not part of the main sample; three are used
as controls and the fourth is excluded for lack of a usable class layer, as described
in Section~\ref{sec:protocol}.

\begin{table}[t]
\centering
\caption{The five audited libraries after manifest construction and recording-level splitting,
i.e.\ the sample every reported number is computed on. Classes are the source classes that
survive the minimum-recording rule; $\bar K$ is the mean number of slices per recording;
$Q_{\text{test}}$ is the sealed panel; $\alpha$ is the realised exposure ratio
$Q^{\text{pool}}_c/m_c$, summarised by its median over classes (not the pooled ratio).
Counts are those of the frozen partition files, so a take released twice under two technique
names is counted once.}
\label{tab:panels}
\footnotesize
\setlength{\tabcolsep}{3.4pt}
\begin{tabular}{lrrrrrrr}
\toprule
Library & Classes & Recordings & Slices & $\bar K$ & $Q_{\text{test}}$ sl. & $Q_{\text{test}}$ rec. & $\alpha$ \\
\midrule
CTIS            & 201 & 4{,}887 & 42{,}505 &  8.70 & 4{,}605 & 963 & 0.160 \\
music genre   &  16 & 1{,}713 & 36{,}375 & 21.23 & 3{,}648 & 337 & 0.166 \\
pianos          &   8 &   668 & 18{,}352 & 27.47 & 1{,}884 & 135 & 0.170 \\
CNPM            &   5 &   286 &  2{,}713 &  9.49 &   271 &  54 & 0.161 \\
bel canto     &   4 &   202 &  9{,}886 & 48.94 &   979 &  40 & 0.161 \\
\midrule
Total           & 234 & 7{,}756 & 109{,}831 & 14.16 & 11{,}387 & 1{,}529 & --- \\
\bottomrule
\end{tabular}
\end{table}

\subsection{Recording identity}

Every slice is mapped to a \texttt{recording\_uid} by a rule specific to the library, derived
from the release format: CTIS by the
per-class file stem, prefixed by the class code when a stem repeats across classes;
music genre by the pair of track index and duration that the file name encodes; pianos by brand,
recording number and duration; CNPM by joining the file name to the dataset's own
\texttt{label.csv}; and bel canto by class and recording number. The identity field parsed for
all 110{,}415 slices; the separate instrument-name field parses
for 99.49\% of CTIS slices, and the 14 unparsed stems are reported in the manifest summary.
We obtained the audio archives of CTIS, CNPM and pianos. Audio-level verification of these
identities was possible only for CTIS, the one library whose audio file names map unambiguously
onto its evaluation recordings; in CNPM and pianos the audio file names carry no reliable link
to the evaluation recording identifiers, so their identities rest on the file-name rule alone.
Bel canto and music genre also release audio, which we did not obtain. A correct parse is in any case not the same as a correct identity: the rule can
reproduce a plausible recording identifier without proving that the released images and the
released audio belong together.

Two checks protect the mapping. First, the mapping between a CTIS recording and the numbering of
its slices is a global counter over classes, not a per-class index; using the per-class index
mis-assigns 218 of 219 classes, and the manifest builder recomputes the global counter before
joining. Second, a SHA-1 sweep over all 4{,}956 CTIS files, 13.3 GB, found two identical pairs --- one take
released under two technique names, and one take with a duplicate-download suffix ---
which we fold to a single canonical identifier carrying one label, so that no canonical
recording appears in two classes and none can sit on two sides of a split by accident. At the slice level CTIS holds 77 groups of identical images,
covering 175 slices. In every one of them the same window appears under more than one recording
name: in 51 groups those names resolve to one recording's audio after the SHA-1 folding above, and
in 26 they resolve to two or more distinct recordings, so the group key rather than the recording
name is what keeps the copies on one side. Separately, in music genre one group of identical
images spans three classes. Within the sealed panel and the exposure arms, every such group is
forced to a single side.

\subsection{Capacity for a recording-level split}\label{sec:capacity}

Capacity planning targets $q_c=\min\{\max(5,\mathrm{round}(0.2n_c)),n_c-1\}$
recordings and $m_c^{\text{plan}}=0.60S_c$ slices. The implemented split also reserves R:
it starts from $\max(1,\mathrm{round}(0.15n_c))$ validation recordings, selected to
approximate $0.15S_c$ slices while retaining the training budget. It starts Q from
$\min\{\max(5,\mathrm{round}(0.2n_c)),n_c-2\}$ recordings, reduces Q when necessary to
leave C non-empty and large enough, and flags these relaxations per class.
The planning half-pool approximation is $\sum_{i\in Q}\lfloor k_i/2\rfloor$;
it is not the exact implemented allocation. Within each selected Q recording, adjacent slice
pairs are joined, and groups sharing an image content hash within the same class are joined
further. Every group is keyed by the smallest content-derived identifier among its members --- a
string formed from the first twelve hexadecimal digits of the slice's SHA-1 --- so the key
depends on image content only, never on time, class or label. Groups are sorted by that key and
assigned alternately to sealed test and injection pool, with singleton groups forced to the test
side and any further copy of an already assigned hash following it. This is neither a contiguous
time half nor independent random assignment; group ordering and duplicate constraints determine
realised pool size.
The realised
exposure ratio $\alpha_c = Q^{\text{pool}}_c/m_c$ is largely determined by design:
a 20\% recording test share, approximately half of its slices in the pool, and a
$0.60S_c$ training budget give $\alpha\approx0.2\times0.5/0.6=0.167$.
Rounding, class support, slice multiplicity and duplicate constraints cause deviations; its median lies between 0.160 and 0.170 in every library, and its per-library span
is 0.000--0.250 for CTIS, 0.162--0.168 for music genre, 0.161--0.173 for pianos, 0.154--0.175
for CNPM and 0.158--0.170 for bel canto. Earlier 50\% and 30\% replacement targets are
capacity questions, distinct from realised exposure. A planning calculation using the upper
95th-percentile pool over repeated draws gives 24 and 54 feasible classes among 219 source
CTIS classes; restricted to the 201 analysed classes these counts are 10 and 38.
These upper-capacity counts do not mean that the frozen split realises such exposures.
CNPM selects 57 recordings for Q, but only 54 contribute sealed test slices: the other three
have all their selected groups allocated to the pool. Table~\ref{tab:panels} counts
recordings actually represented in the test panel.

Two consequences shape the design. First, equal replacement, not addition: injecting
$Q^{\text{pool}}_c$ slices while keeping $m_c$ fixed removes exactly as many C slices as it adds,
so the arms differ in composition and not in size. Second, the training-recording count must be
matched on the C side: A0 and A1 draw their training recordings from C, and matching the total
training-recording count across arms is impossible once A1 must also include pool recordings, so
we match the C-side count exactly and report the A1 total, C side plus pool, per class. Where
the C side cannot supply the planned budget the realised $m_c$ is capped at what C holds, and
where a class cannot yield the planned $q_c$ the pool is trimmed; both events are flagged per
class in the released feasibility table, so a reader can see exactly where the arithmetic binds.

We required at least five recordings per class as a study-specific eligibility rule; this is not a
mathematical minimum for every possible recording-level split. Applying that rule to CTIS reduces 219 source
classes to 206. Five further classes (D0051, D0058, D0060, D0061 and D0062) each contain
one placeholder recording whose identifier includes a Chinese ``no audio'' marker; removing it reduces
five recordings to four in each class, giving 201 classes and 4{,}887 recordings. The placeholder
slices in D0060--D0062 share a byte-identical set of 21 images; D0051 and D0058 have 27 and 26
placeholder images, respectively. Supplementary Material S1 lists the class-level exclusions. The other four libraries keep
all their classes; two of them lose one recording each, because duplicate removal took every
slice of that recording, leaving CNPM with 286 of 287 and bel canto with 202 of 203, while
pianos and music genre lose nothing. We report both counts wherever the source documents state 219, so that the
difference between the released class count and the auditable one is visible rather than
implicit.

\subsection{Feature cache}

All five libraries release MEL images at 496$\times$369. We resize to 224$\times$224 with
bilinear interpolation and store the result as a \texttt{uint8} memmap in manifest order: the
109{,}831 audited images plus the excluded classes come to 15.48 GiB for the five libraries, and
the 5{,}432 negative-control images add 0.76 GiB. Aspect ratio is deliberately not preserved: the source
images share one size here, but across the full release the aspect ratios span two orders of
magnitude, and padding to a common shape would turn padding fraction into a library-level
covariate. A random sample of 200 images per library was re-decoded and compared with the cache
byte for byte: the maximum absolute difference is zero for all eight libraries.
